\section{Vaught's never-two theorem and countable-model problems}
\label{sec:p03-vaught-never-two-presentation}

Vaught's original Section~6 answers a question of Raphael Robinson. The
introduction of the paper announces that Section~6 shows that a complete
theory cannot have exactly two nonisomorphic denumerable models, and that
this answers Robinson's question \cite[p.~304]{Vaught1961}. Section~6 then
states the question. Robinson had observed that the complete theory of an
equivalence relation with exactly two classes, both infinite, has exactly two
nonisomorphic models of power \(\aleph_1\), and he asked whether some
complete theory has exactly two nonisomorphic models of power \(\aleph_0\).
A footnote dates the question to a conversation in 1957 among several
people, Vaught included \cite[p.~318 and n.~23]{Vaught1961}. Vaught's
answer combines examples with a restriction: Theorem~6.1 says that a
complete theory with exactly \(n\) nonisomorphic denumerable models exists
if and only if \(n\ne2\) \cite[pp.~318--320]{Vaught1961}. The paper's letters
\(i,\ldots,r\) range over the natural numbers \(0,1,\ldots\)
\cite[p.~304, n.~5]{Vaught1961}, so the theorem covers every \(n\in\omega\),
including \(n=0\); Theorem~\ref{thm:p03-w03-p03-016-theorem-6-1} states
it in that form. In Vaught's usage a complete theory is consistent
\cite[p.~303, n.~1]{Vaught1961}. Throughout this presentation a theory is
consistent, and countable includes finite unless ``denumerable'' is specified.
All languages here are countable. Until the final graph translation in the
never-two proof, the language has only relation symbols and constants, as in
the original presentation.
The graph translation described in the proof transfers the negative result
to languages with function symbols.

\subsection{Type counting and the named-type theory}
\label{sec:p03-vaught-never-two-types}

\begin{hbdefinition}{def:p03-vaught-spectrum-notation}{Spectra, formula algebras and named types}
For a complete first-order theory \(T\), \(I(T,\aleph_0)\) counts its
countable models up to isomorphism. Write \(S_m(T)\) for its complete
\(m\)-types and \(F_m(T)\) for formulas in the first \(m\) free variables
modulo equivalence in \(T\). Boolean conjunction, disjunction and negation
make \(F_m(T)\) a Boolean algebra. Vaught's \(\mathfrak P_m(T)\), the set
of dual prime ideals of this algebra, corresponds to \(S_m(T)\): a type
contains exactly the formulas in its associated dual prime ideal.

For \(p\in S_m(T)\) and fresh constants \(\bar c=(c_0,\ldots,c_{m-1})\),
put \(T_p=T\cup\{\varphi(\bar c):\varphi\in p\}\).
This is a complete theory in the expanded language. Its models are precisely
pointed models \((M,\bar a)\) with \(\operatorname{tp}_M(\bar a)=p\).
Inclusion of the original language induces a map from \(F_k(T)\) to
\(F_k(T_p)\); injectivity is a theorem below, not part of this definition.
\end{hbdefinition}

The complete-type and isolation conventions are those of
Definition~\ref{def:principal-type}. Saturation here concerns finite tuples
over finite parameter sets, as in Theorem~\ref{thm:p03-countable-saturated}.
The original existence and back-and-forth constructions are supplied in
Theorems~\ref{thm:p03-vaught-saturated-existence} and
\ref{thm:p03-vaught-saturated-embedding}.

\begin{hblemma}{thm:p03-vaught-few-models-types}{proved=annex:p03.vaught-two.types}{Few countable models give small type spaces}
If \(T\) has an infinite model and \(I(T,\aleph_0)\leq\aleph_0\), every
\(S_m(T)\) is countable. In every countable type space, isolated types are
dense. A type space in which every type is isolated is finite; conversely,
every type in a finite type space is isolated.
\end{hblemma}

Thus the hypothesis of exactly two models supplies smallness. It does not
supply \(\omega\)-stability, which counts types over all countable parameter
sets. The original prime/atomic construction
(Theorem~\ref{thm:p03-vaught-prime-atomic}) and saturation construction
operate under their displayed hypotheses; neither adjective may be exchanged
for the other in this argument.

\begin{hblemma}{thm:p03-vaught-named-finiteness}{proved=annex:p03.vaught-two.named}{Categoricity after naming a complete type}
Let \(T\) be complete with infinite models. If \(T_p\) has exactly one
denumerable model up to isomorphism, every \(F_k(T_p)\) is finite.
Moreover, \(F_k(T)\to F_k(T_p)\) is injective for every \(k\).
Consequently, if some \(S_m(T)\) contains a nonisolated type, \(T_p\)
cannot be denumerably categorical for any \(p\).
\end{hblemma}

\begin{hbtheorem}{thm:p03-vaught-never-two-original}{proved=annex:p03.vaught-two.never}{Vaught's never-two theorem}
A complete first-order theory in a countable language cannot have exactly
two nonisomorphic countable models. In particular,
\(I(T,\aleph_0)\ne2\).
\end{hbtheorem}

The proof names a realization of a nonisolated type. A model omitting that
type is one of the hypothetical two isomorphism types. Every model realizing
it must therefore be the other, saturated type. Back-and-forth then makes
\(T_p\) categorical, contradicting
Lemma~\ref{thm:p03-vaught-named-finiteness}. Finite models and the absence
of a nonisolated type are separate cases in the annex.

\subsection{The examples surrounding the theorem}
\label{sec:p03-vaught-finite-count-examples}

\begin{hbdefinition}{def:p03-vaught-colored-orders}{Ehrenfeucht's orders and their finite-language coding}
Fix a positive integer \(r\). Let \(T_{r+2}\) assert that \(<\) is a dense
linear order without endpoints, the unary predicates
\(U_0,\ldots,U_{r-1}\) partition its domain and are each dense, and
\(c_0<c_1<\cdots\), with every \(c_i\in U_0\).
For \(r=1\), discard the redundant predicate and obtain the original
three-model theory \(T_3\).

Its finite-language coding \(T'_{r+2}\) has an equivalence relation \(R\),
an order \(<\) constant on \(R\)-classes, and class-constant predicates
\(U_j\). On the quotient the order is dense without endpoints and each
color is dense. For every positive integer \(i\), exactly one class
\(C_i\) has size \(i\); these classes have color \(U_0\) and satisfy
\(C_i<C_{i+1}\). All other classes are infinite. ``Class has size \(i\)''
is its ordinary first-order formula with \(i\) distinct members and a
universal assertion that these exhaust the class. The notation \(C_i\)
denotes a definable quotient point, not an added constant symbol.
For \(r=1\), this is Vaught's displayed language \((<,R)\).
\end{hbdefinition}

\begin{hbtheorem}{thm:p03-vaught-finite-counts}{proved=annex:p03.vaught-two.examples}{All positive finite counts except two}
For every positive integer \(n\), a complete first-order theory has exactly
\(n\) nonisomorphic denumerable models if and only if \(n\ne2\).
For \(n=1\), dense linear order without endpoints suffices. For \(n=r+2\),
both \(T_{r+2}\) and \(T'_{r+2}\) are complete and have exactly \(r+2\)
denumerable models. The latter has a finite language.
\end{hbtheorem}

For the named-constant theory the possibilities are: the constants are
cofinal; they have upper bounds but no least upper bound in the domain; or
their least upper bound exists and has one of the \(r\) colors. The quotient
points \(C_1,C_2,\ldots\) give the same alternatives for the coding.
The upper-bound condition is the infinite conjunction
\(\bigwedge_i c_i<x\); it is not an additional first-order predicate.
The annex proves completeness as well as the isomorphism classification.

Vaught's own justifications in this part of the proof are short citations.
That the theory \(T_1\) of dense linear order without endpoints is complete
with exactly one denumerable model up to isomorphism is credited to Cantor's
theorem \cite[p.~318]{Vaught1961}. For \(T_n\) with \(n\ge4\), which is
\(T_{r+2}\) above with \(r=n-2\), he invokes a theorem of
Skolem (his reference~[14]): any two denumerable models of \(T_n\) are
isomorphic once their lists of distinguished elements are ignored. He states
that completeness and the list of isomorphism types are easily seen from this
fact \cite[p.~319]{Vaught1961}. The colored back-and-forth argument at the
start of the annex proof of Theorem~\ref{thm:p03-vaught-finite-counts} is a
proof of Skolem's fact as Vaught uses it, and its uncolored case is Cantor's
theorem. Clause~1 of Theorem~\ref{thm:p03-w03-p03-016-source-routes}
derives the completeness of \(T_{r+2}\) from Skolem's fact in the way the
source indicates.

For the finite-language theory Vaught sketches a different completeness
argument, granted that the theory has only three denumerable models
\cite[pp.~319--320]{Vaught1961}. Start from a denumerable model of \(T'_3\)
in which the finite classes are cofinal. Adding a constant above all of
them and applying the completeness theorem yields an elementarily equivalent
denumerable model of one of the other two kinds. Further applications of the
completeness theorem, together with the union of an elementarily increasing
family, show that a model in which the finite classes have a limit is
elementarily equivalent to one in which they neither have a limit nor are
cofinal. Completeness then follows from the L\"owenheim--Skolem theorem.
Clause~2 of Theorem~\ref{thm:p03-w03-p03-016-source-routes} carries out this
route in full for every \(T'_{r+2}\); the quantifier-rank game in the annex
proof of Theorem~\ref{thm:p03-vaught-finite-counts} is an independent second
proof.

\begin{hbtheorem}{thm:p03-w03-p03-016-source-routes}{proved=annex:p03-w03-p03-016-source-routes}{Completeness by Vaught's two routes}
Let \(r\) be a positive integer.
\begin{enumerate}
  \item (Skolem's fact.) Every sentence of the language of \(T_{r+2}\)
        mentions only \(c_0,\ldots,c_k\) for some \(k\). The reducts of any
        two denumerable models of \(T_{r+2}\) to
        \(<,U_0,\ldots,U_{r-1},c_0,\ldots,c_k\) are isomorphic. Hence
        \(T_{r+2}\) is complete.
  \item (The finite-language sketch.) Let \(A\) be a denumerable model of
        \(T'_{r+2}\) in which the finite classes are cofinal. Some
        denumerable model of \(T'_{r+2}\) in which the finite classes are
        bounded is elementarily equivalent to \(A\). Every denumerable model
        of \(T'_{r+2}\) in which the finite classes have a least upper bound
        is elementarily equivalent to a denumerable model in which they are
        bounded and have no least upper bound. Consequently all \(r+2\)
        denumerable isomorphism types are elementarily equivalent, and
        \(T'_{r+2}\) is complete.
\end{enumerate}
\end{hbtheorem}

\begin{hbtheorem}{thm:p03-w03-p03-016-theorem-6-1}{proved=annex:p03-w03-p03-016-theorem-6-1}{Theorem~6.1 over all natural numbers}
For every \(n\in\omega\) there is a complete theory, in a countable language
whose nonlogical symbols are relation symbols and individual constants, with
exactly \(n\) nonisomorphic denumerable models, if and only if \(n\ne2\).
For \(n=0\) the complete theory of a one-element structure is a witness.
Every complete theory in a countable language with an infinite model has at
least one denumerable model.
\end{hbtheorem}

The case \(n=0\) is not discussed in Vaught's text. Its witnesses have only
finite models, so they lie outside the standing assumption, adopted in
\S\,1, that \(T\) is a complete theory with infinite models
\cite[p.~306]{Vaught1961}. The last sentence of the theorem shows that,
inside that assumption, the finite counts which occur are exactly the
positive \(n\ne2\).

\begin{hbtheorem}{thm:p03-vaught-equivalence-example}{proved=annex:p03.vaught-two.equivalence}{Two models at an uncountable cardinal}
The theory asserting exactly two infinite equivalence classes is complete.
It has one denumerable model, but exactly two models of size \(\aleph_1\):
the unordered class-size pairs are
\(\{\aleph_0,\aleph_1\}\) and \(\{\aleph_1,\aleph_1\}\).
\end{hbtheorem}

This is the theory from which Robinson's question starts
\cite[p.~318]{Vaught1961}. Its two models occur at the cardinal
\(\aleph_1\), not at the countable cardinal. The never-two theorem answers
the question negatively: no complete theory has exactly two nonisomorphic
denumerable models. If finite models are counted as countable, a complete
theory with a finite model has exactly one countable isomorphism type, and a
complete theory with infinite models has a denumerable one, so that count is
never \(0\). Counting denumerable models, as Vaught does, the value \(0\)
occurs, by Theorem~\ref{thm:p03-w03-p03-016-theorem-6-1}.

\subsection{The original questions and the later formulations}
\label{sec:p03-vaught-original-problem}

Immediately after the proof of Theorem~6.1, Vaught records a conjecture
which he calls unresolved and which, he says, appears to have been made by
a number of people. In his words: ``A theory which is
\(\aleph_{1+\alpha}\)-categorical but not \(\aleph_0\)-categorical has
exactly \(\aleph_0\) non-isomorphic denumerable models.''
\cite[p.~320]{Vaught1961}. The count is of denumerable models, so finite
models play no part in it, and the hypothesis is categoricity in some
uncountable power \(\aleph_{1+\alpha}\). Theories here have at most
\(\aleph_0\) nonlogical constants, each a relation symbol or an individual
constant \cite[p.~304]{Vaught1961}. From a point in \S\,1 onward the paper
assumes that the theory \(T\) under discussion is complete with infinite
models, while remarking that one or both assumptions can often be dropped
\cite[p.~306]{Vaught1961}. The conjecture itself says only ``a theory''.
He then poses a second problem: whether it can be proved, without the
continuum hypothesis, that some complete theory has exactly \(\aleph_1\)
nonisomorphic denumerable models \cite[p.~320]{Vaught1961}. These are the
actual original questions. The cardinal dichotomy customarily called
Vaught's conjecture is formulated below; it does not appear in this form in
the 1961 text. The introduction of the selected 2026 status source
attributes the countable-or-continuum statement to this paper
\cite[p.~1]{GonzalezRosseggerTuretsky2026Status}; the 1961 text itself
displays only the two questions just quoted.

The categoricity conjecture is a statement about uncountably categorical
theories, and it is distinct from the never-two theorem: the never-two
theorem excludes one finite count for every complete theory and says nothing
about which theories have \(\aleph_0\) denumerable models. For a complete
theory in a countable language, the later Baldwin--Lachlan refinement
(Theorem~\ref{thm:p03-baldwin-lachlan}, a cited result) gives
\(I(T,\aleph_0)\in\{1,\aleph_0\}\) under categoricity in one uncountable
cardinal. Such a theory has no finite models, so this count is a count of
denumerable models. If \(T\) is not \(\aleph_0\)-categorical, the value
\(1\) is excluded and the count is \(\aleph_0\), which is Vaught's
conjecture for complete theories.

\begin{hbdefinition}{def:p03-vaught-variants}{Cardinal, perfect-set and infinitary claims}
The \emph{first-order cardinal conjecture} says, for every complete theory
in a countable language,
\[
 I(T,\aleph_0)\leq\aleph_0\quad\text{or}\quad
 I(T,\aleph_0)=2^{\aleph_0}.
\]
The \emph{perfect-set formulation} asks instead for a perfect set of
pairwise nonisomorphic models whenever there are uncountably many
isomorphism types. Models with domain \(\omega\) are coded by the product
space of their relation, function and constant interpretations; finite
domains have the analogous coding components. Perfect means nonempty,
closed and without isolated points in the resulting coding space.

The \emph{infinitary cardinal conjecture} makes the same counting assertion
for the countable models of each \(L_{\omega_1\omega}\)-sentence
\(\sigma\). This logic permits countable conjunctions and disjunctions,
with only finitely many free variables in such a connective and finite
quantifier strings. It is a wider axiomatization class than a complete
first-order theory. The perfect-set question can likewise be asked in this
infinitary class.
\end{hbdefinition}

\begin{hblemma}{thm:p03-vaught-cardinal-boundaries}{proved=annex:p03.vaught-two.variants}{What cardinal counting does and does not assert}
In a countable language there are at most \(2^{\aleph_0}\) countable
isomorphism types. A perfect set of pairwise nonisomorphic models supplies
\(2^{\aleph_0}\) types. Under CH, the bare countable-or-continuum cardinal
dichotomy follows from this upper bound for any class of countable
structures. That counting argument supplies no perfect set.
There is an \(L_{\omega_1\omega}\)-sentence with exactly two denumerable
models whose first-order reducts both satisfy the complete theory \(T_3\).
\end{hblemma}

Every countable first-order theory can be conjoined into one infinitary
sentence, so the infinitary cardinal claim entails the first-order one.
The reverse implication is not supplied by this translation. Likewise,
cardinality alone does not supply the topological witnesses in the perfect-set
formulation. Morley's separate 1970 source treatment concerns its precise
infinitary threshold; the never-two proof does not establish that theorem.

\subsection{Scott ranks and the dated status}
\label{sec:p03-vaught-status-2026}

\begin{hbdefinition}{def:p03-vaught-scott-comparison}{Ranks used in the status comparison}
A Scott sentence of a countable structure characterizes it up to isomorphism
among countable structures. In the selected 2026 convention,
\(SR(M)\) is the least \(\alpha\) for which such a sentence is
\(\Pi^{\mathrm{in}}_{\alpha+1}\), and \(SR_p(M)\) is the least
\(\alpha\) for which one is \(\Sigma^{\mathrm{in}}_{\alpha+2}\).
Here \(\Sigma^{\mathrm{in}}_0=\Pi^{\mathrm{in}}_0\) are finitary
quantifier-free formulas; higher \(\Sigma\)-classes are countable
disjunctions of existentially quantified lower \(\Pi\)-formulas, and the
\(\Pi\)-classes are the dual conjunctions of universally quantified lower
\(\Sigma\)-formulas. Parameterized rank permits one fixed finite tuple of
named parameters; it is not interchangeable with unparameterized rank.

Write \(\operatorname{vo}(\sigma)\) for the Vaught ordinal: the least
ordinal \(\beta\), if one exists, such that either continuum many
\(\Pi^{\mathrm{in}}_\beta\)-types are realized among the models of
\(\sigma\), or there are only countably many models of \(\sigma\) and they
all have Scott rank less than \(\beta\); otherwise use \(\omega_1\)
\cite[p.~6]{GonzalezRosseggerTuretsky2026Status}. Here the
\(\Pi^{\mathrm{in}}_\beta\)-type of a tuple is the set of
\(\Pi^{\mathrm{in}}_\beta\)-formulas it satisfies, and, as Scott ranks
require, the models are the countable models of \(\sigma\). Since negation
exchanges \(\Pi^{\mathrm{in}}_\beta\) and \(\Sigma^{\mathrm{in}}_\beta\)
up to logical equivalence, the \(\Pi^{\mathrm{in}}_\beta\)-type and the
\(\Sigma^{\mathrm{in}}_\beta\)-type of a tuple determine each other.

The \emph{\(\omega\)-Vaught conjecture}, in the formulation the selected
source attributes to Gonzalez and Montalb\'an, states that every theory
has a countable Vaught ordinal and, more strongly, that every
\(\Pi^{\mathrm{in}}_\alpha\)-theory \(\sigma\) satisfies
\(\operatorname{vo}(\sigma)\leq\alpha+\omega\)
\cite[p.~4]{GonzalezRosseggerTuretsky2026Status}. In the infinitary setting
the source describes it as the translation of Martin's conjecture, a
model-theoretic strengthening of Vaught's conjecture
\cite[p.~4]{GonzalezRosseggerTuretsky2026Status}.
\end{hbdefinition}

The selected source states that the existence of a countable Vaught ordinal
for \(\sigma\) is equivalent to \(\sigma\) satisfying Vaught's conjecture,
in its infinitary formulation, and that a counterexample has Vaught ordinal
\(\omega_1\) \cite[pp.~2, 6]{GonzalezRosseggerTuretsky2026Status}. The
\(\omega\)-Vaught conjecture therefore strengthens the infinitary conjecture:
it keeps the countability of every Vaught ordinal and adds the bound
\(\alpha+\omega\) determined by the syntactic complexity of the sentence.
Theorem~\ref{thm:p03-w03-p03-016-vo-cardinal} proves, in ZFC, the direction
from a countable Vaught ordinal to the cardinal dichotomy, and explains why
the reverse direction cannot be read as a statement about the bare cardinal
dichotomy alone once the continuum hypothesis is assumed.

\paragraph{Edition status, 9 September 2026.}
The selected evidence is Gonzalez, Rossegger and Turetsky,
arXiv:2606.15205v2, dated 26 August 2026. Its introduction treats the general
conjecture as open and subsequently uses the infinitary formulation.
Corollary~2.11 concerns hypothetical \(\Pi^{\mathrm{in}}_2\)
counterexamples: they must have two models at every positive countable
parameterized Scott rank. Theorem~3.6 gives a different conclusion: a minimal
counterexample has exactly one model at each rank in an unbounded set of
unparameterized Scott ranks. These statements do not contradict one another.
Corollary~4.10 proves the \(\omega\)-Vaught bound for infinitary theories
with finitely many models, not for arbitrary theories. This restricted
result and the never-two theorem leave the general conjecture unresolved
\cite{GonzalezRosseggerTuretsky2026Status}.

The version record reports added assumptions in Section~3. The displayed
status therefore belongs to v2 and to the edition's historical date; it does
not silently combine v1 statements, later versions, or the three different
counting questions.
